% Einheit 2 von 2 – Diagramme beurteilen und Boxplot (Katalog: daten.md, Einheit 5)
\setcounter{aufgabe}{23}
\zweigkopf{e2}{Einheit 2 von 2 · Diagramme beurteilen und Boxplot}{Hier lernst du, Aussagen zu Diagrammen zu prüfen, täuschende Diagramme zu erkennen und Boxplots zu lesen · neu oder schon bekannt – je nach Buch, bis Klasse 9 (am Gymnasium neu in diesem Jahr) · P10 · baut auf: Säulendiagramme lesen und Kenngrößen (Kennst du schon), Prozentsatz, Bruchteil und Vielfaches}

\begin{aufgabe}{\textbf{Ich erkenne, wo die senkrechte Achse eines Diagramms anfängt.} Kreuze an, ob die Achse bei 0 beginnt. Wenn nicht, schreibe den Startwert auf. Lies keine Säule ab.}
\noindent\begin{minipage}[t]{0.48\linewidth}
\begin{teile}
\teil Beginnt bei 0? \janein\\ sonst Startwert: \leerfeld
\end{teile}
\saeulenab[ymax=40,ystep=10,yfein=5,ylabel=Anzahl]{A/25,B/35,C/15}
\end{minipage}\hfill
\begin{minipage}[t]{0.48\linewidth}
\begin{teile}
\teil Beginnt bei 0? \janein\\ sonst Startwert: \leerfeld
\end{teile}
\saeulenab[ymin=80,ymax=100,ystep=5,yfein=1,ylabel=Punkte]{A/92,B/85,C/97}
\end{minipage}

\noindent\begin{minipage}[t]{0.48\linewidth}
\begin{teile}
\teil Beginnt bei 0? \janein\\ sonst Startwert: \leerfeld
\end{teile}
\saeulenab[ymin=1.5,ymax=1.9,ystep=0.1,yfein=0.05,ylabel=Preis in €]{Mo/1.75,Di/1.6,Mi/1.85}
\end{minipage}\hfill
\begin{minipage}[t]{0.48\linewidth}
\begin{teile}
\teil Beginnt bei 0? \janein\\ sonst Startwert: \leerfeld
\end{teile}
\saeulenab[ymax=500,ystep=100,yfein=50,ylabel=Besucher]{Sa/450,So/300,Mo/150}
\end{minipage}
\end{aufgabe}

\begin{aufgabe}{\textbf{Ich kann prüfen, ob eine Aussage zu einem Diagramm stimmt.} Das Diagramm zeigt die Besucher des Freibads „Sonnenwiese“ im Sommer. Kreuze an, ob die Aussage wahr oder falsch ist.}

\saeulenab[ymax=50,ystep=10,yfein=2,ylabel=Besucher in Tausend]{Mai/18,Juni/34,Juli/46,August/41,September/12}

\begin{teile}
\teil Im Juli kamen die meisten Besucher. \hfill \kreuz{wahr}\quad\kreuz{falsch}
\teil Im Mai kamen 18\,000 Besucher. \hfill \kreuz{wahr}\quad\kreuz{falsch}
\teil Im September kamen mehr als 15\,000 Besucher. \hfill \kreuz{wahr}\quad\kreuz{falsch}
\teil Im August kamen 4100 Besucher. \hfill \kreuz{wahr}\quad\kreuz{falsch}
\teil Von Juni bis August stieg die Zahl der Besucher jeden Monat. \hfill \kreuz{wahr}\quad\kreuz{falsch}
\teil Im Juli kamen 12\,000 Besucher mehr als im Juni. \hfill \kreuz{wahr}\quad\kreuz{falsch}
\end{teile}
\end{aufgabe}

\begin{aufgabe}{\textbf{Ich kann prüfen, ob eine Aussage zu einem Diagramm stimmt – weiter.} 600 Jugendliche wurden gefragt, was sie in der Freizeit am liebsten machen; jeder durfte eine Antwort geben. Kreuze an und begründe mit einer Rechnung.}

\balkenab[xmax=35,xstep=5,xfein=1,xlabel=Anteil in \%]{Freunde treffen/32,Gaming/20,Sport/24,Musik hören/15,Lesen/9}

\begin{teile}
\teil Mehr als die Hälfte trifft am liebsten Freunde oder spielt am liebsten Games. \hfill \kreuz{wahr}\quad\kreuz{falsch}
\teil Gaming wird um ein Drittel häufiger genannt als Musik hören. \hfill \kreuz{wahr}\quad\kreuz{falsch}
\teil Jede 15. befragte Person hört am liebsten Musik. \hfill \kreuz{wahr}\quad\kreuz{falsch}
\teil Sport ist beliebter als Gaming, und mehr als 60 Jugendliche lesen am liebsten. \hfill \kreuz{wahr}\quad\kreuz{falsch}
\teil Kreuze zu jeder Aussage an und begründe. (P10 2018)\\
(1) Weniger als die Hälfte trifft am liebsten Freunde oder liest am liebsten.\\ \kreuz{wahr}\quad\kreuz{falsch}\quad\kreuz{nicht entscheidbar}\\
(2) Jede 5. befragte Person macht am liebsten Sport.\\ \kreuz{wahr}\quad\kreuz{falsch}\quad\kreuz{nicht entscheidbar}\\
(3) Mehr Mädchen als Jungen lesen am liebsten.\\ \kreuz{wahr}\quad\kreuz{falsch}\quad\kreuz{nicht entscheidbar}
\end{teile}
\end{aufgabe}

\begin{aufgabe}{\textbf{Ich kann erklären, warum ein Diagramm täuscht.} Ein Turnverein zeigt seine Mitgliederzahlen so:}

\saeulenab[ymin=400,ymax=450,ystep=10,yfein=2,ylabel=Mitglieder]{2022/410,2023/420,2024/440}

\begin{teile}
\teil Bei welchem Wert beginnt die senkrechte Achse? \leerfeld
\teil Die Säule für 2024 ist doppelt so hoch wie die für 2023. Hat sich die Zahl der Mitglieder verdoppelt? Begründe mit den Zahlen.
\teil Um wie viel Prozent ist die Zahl der Mitglieder von 2023 bis 2024 wirklich gestiegen? Runde auf eine Nachkommastelle. \leerfeld[\%]
\teil Der Vorstand zeichnet dieselben Zahlen noch einmal, aber mit 1\,cm für 5 Mitglieder statt 1\,cm für 10 Mitglieder. Wie ändert sich der Eindruck?
\end{teile}
Ein Obsthändler zeigt dieses Diagramm und sagt: „Die Apfelpreise sind explodiert!“

\saeulenab[ymin=2.3,ymax=2.7,ystep=0.1,yfein=0.05,ylabel=Preis in € je kg]{W1/2.4,W2/2.45,W3/2.5,W4/2.55,W5/2.6}

\begin{teile}
\teil Erkläre, warum das Diagramm einen falschen Eindruck erweckt. Gib an, um wie viel Prozent der Preis von der ersten Woche (W1) bis zur fünften Woche (W5) wirklich gestiegen ist. (P10 2026)
\end{teile}
\end{aufgabe}

\begin{aufgabe}{\textbf{Ich kann einen Trend in einem Diagramm beurteilen (kommt in Klasse 9, am Gymnasium neu in diesem Jahr).} Ein Kiosk hat notiert, wie viele Zeitungen er im Durchschnitt an einem Tag verkauft.}

\liniendia[ymax=140,ystep=20,yfein=5,ylabel=Zeitungen je Tag]{2019/120,2020/105,2021/95,2022/100,2023/80,2024/70,2025/60}

\begin{teile}
\teil Beschreibe den Verlauf von 2019 bis 2025 in einem Satz.
\teil „Die Zahl der verkauften Zeitungen ist von Jahr zu Jahr gesunken.“ Stimmt das? Begründe.
\teil Der Kioskbesitzer sagt: „Wenn es so weitergeht, verkaufe ich 2031 keine einzige Zeitung mehr.“ Wie kommt er auf 2031? Warum ist die Aussage trotzdem unsicher?
\end{teile}
\end{aufgabe}

\begin{aufgabe}{\textbf{Ich kann ein Diagramm mit einer Achse ab 0 zeichnen und eine Zeitungsaussage dazu prüfen.} Das linke Diagramm zeigt, wie viele Bücher die Stadtbibliothek ausgeliehen hat.}

\noindent\begin{minipage}[t]{0.48\linewidth}
\saeulenab[ymin=30,ymax=45,ystep=5,yfein=1,breite=0.7,ylabel=Ausleihen in Tausend]{2021/42,2022/40,2023/41,2024/37,2025/35}
\end{minipage}\hfill
\begin{minipage}[t]{0.48\linewidth}
\saeulenab[ymax=50,ystep=10,yfein=2,breite=0.7,ylabel=Ausleihen in Tausend]{2021/,2022/,2023/,2024/,2025/}
\end{minipage}

\begin{teile}
\teil Lies die Werte ab und trage sie ein.
\end{teile}
\sachtabelle{lccccc}{Jahr & 2021 & 2022 & 2023 & 2024 & 2025}{Ausleihen in Tausend & \leerzelle & \leerzelle & \leerzelle & \leerzelle & \leerzelle}
\begin{teile}
\teil Zeichne die Säulen in das rechte Diagramm, dessen Achse bei 0 beginnt.
\teil Vergleiche den Eindruck der beiden Diagramme in einem Satz.
\teil In der Zeitung steht zum linken Diagramm: „Bücher leihen ist out – die Ausleihen sind seit 2021 jedes Jahr gesunken und haben sich fast halbiert.“ Prüfe beide Teile der Aussage mit den Zahlen und nenne, was am linken Diagramm den Eindruck verstärkt. (P10 2019)
\end{teile}
\end{aufgabe}

\begin{aufgabe}{\textbf{Ich kann einen Boxplot lesen (kommt nächstes Jahr, je nach Buch bis Klasse 9; am Gymnasium neu in diesem Jahr).} Ein Boxplot zeigt fünf Werte: den kleinsten Wert, das untere Quartil, den Median, das obere Quartil und den größten Wert. In der Box liegt die mittlere Hälfte der Werte.}
Beispiel: Boxplot mit den Werten $2 \mid 5 \mid 8 \mid 10 \mid 15$: kleinster Wert 2, unteres Quartil 5, Median 8, oberes Quartil 10, größter Wert 15; Spannweite $15 - 2 = 13$.

Die Boxplots zeigen die Punkte im Mathetest der Klassen 7a und 7b.

\begin{boxplots}[xmin=0,xmax=24,xstep=1,xlabel=Punkte]
\bp{6,10,13,16,22}{Klasse 7a}
\bp{9,12,14,15,19}{Klasse 7b}
\end{boxplots}

\begin{teile}
\teil Median der 7a: \leerfeld
\teil kleinster und größter Wert der 7a: \leerfeld \quad \leerfeld
\teil unteres und oberes Quartil der 7a: \leerfeld \quad \leerfeld
\teil Spannweite der 7a: \leerfeld
\teil Zwischen welchen Punktzahlen liegt die mittlere Hälfte der 7b? \leerfeld
\teil Welche Klasse hat den höheren Median? Bei welcher Klasse liegen die Punkte weiter auseinander? Begründe mit zwei Werten aus den Boxplots.
\end{teile}
\end{aufgabe}

\begin{aufgabe}{\textbf{Ich kann einen Boxplot zeichnen (kommt nächstes Jahr, je nach Buch bis Klasse 9; am Gymnasium neu in diesem Jahr).} Zeichne den Boxplot in die Zeile mit dem passenden Buchstaben.}
\begin{teile}
\teil kleinster Wert 3, unteres Quartil 7, Median 10, oberes Quartil 12, größter Wert 18
\teil kleinster Wert 5, unteres Quartil 6, Median 11, oberes Quartil 15, größter Wert 20
\end{teile}
Das untere Quartil ist der Median der unteren Hälfte (ohne den Median), das obere Quartil der Median der oberen Hälfte.
\begin{teile}
\teil Elf Kinder haben beim Zielwerfen diese Treffer: $9;\ 4;\ 13;\ 5;\ 7;\ 11;\ 2;\ 9;\ 5;\ 8;\ 6$. Bestimme die fünf Werte und zeichne den Boxplot.\\ kleinster Wert: \leerfeld \quad unteres Quartil: \leerfeld\\ Median: \leerfeld \quad oberes Quartil: \leerfeld\\ größter Wert: \leerfeld
\end{teile}

\begin{boxplots}[xmin=0,xmax=20,xstep=1,xlabel=Wert]
\bp{}{c)}
\bp{}{b)}
\bp{}{a)}
\end{boxplots}
\end{aufgabe}

\begin{aufgabe}{\textbf{Ich finde den Fehler, wenn jemand nur die Säulenhöhe vergleicht.} Das Diagramm zeigt die Besucher eines Stadtfests.}

\noindent\begin{minipage}[t]{0.48\linewidth}
\saeulenab[ymin=200,ymax=340,ystep=20,yfein=10,ylabel=Besucher]{Samstag/260,Sonntag/320}
\end{minipage}\hfill
\begin{minipage}[t]{0.48\linewidth}
\saeulenab[ymin=50,ymax=70,ystep=5,yfein=1,ylabel=Spenden in €]{Januar/55,Februar/65}
\end{minipage}

\noindent Mo schreibt zum linken Diagramm: „Die Säule für Sonntag ist doppelt so hoch wie die für Samstag. Also kamen am Sonntag doppelt so viele Besucher.“
\begin{teile}
\teil Finde den Fehler, benenne ihn und korrigiere die Aussage. Um wie viel Prozent waren es am Sonntag mehr? \leerfeld[\%]
\teil Das rechte Diagramm zeigt die Spenden einer Klasse. Wie viel Mal so hoch wirkt die Säule für Februar? Um wie viel Prozent sind die Spenden wirklich gestiegen? Runde auf eine Nachkommastelle.\\ wirkt \leerfeld-mal so hoch \quad wirklich gestiegen um \leerfeld[\%]
\end{teile}
\end{aufgabe}

\begin{aufgabe}{\textbf{Ich kann begründen, wann die Achse eines Diagramms bei 0 beginnen sollte.} Antworte in ganzen Sätzen.}
\begin{teile}
\teil Begründe, warum die senkrechte Achse eines Säulendiagramms bei 0 beginnen sollte.
\teil In einem Liniendiagramm zur Körpertemperatur eines kranken Kindes über drei Tage beginnt die Achse bei $35\,°$C. Ist das hier auch eine Täuschung? Begründe.
\end{teile}
\end{aufgabe}

\begin{aufgabe}{\textbf{Ich kann mit Zahlen aus einer Zählung einen Antrag begründen.} Die Schülervertretung hat die Fahrräder auf dem Schulhof gezählt: März 84, April 96, Mai 118, Juni 125. Es gibt 5 Fahrradständer mit je 20 Plätzen.}
\begin{teile}
\teil Wie viele Plätze fehlten im Juni? Wie viele Ständer muss die Schule mindestens dazukaufen?\\ fehlende Plätze: \leerfeld \quad Ständer: \leerfeld
\teil Um wie viel Prozent ist die Zahl der Fahrräder von März bis Juni gestiegen? Runde auf eine Nachkommastelle. \leerfeld[\%]
\teil Die Schülervertretung überlegt, im Antrag ein Säulendiagramm zu zeigen, dessen Achse erst bei 80 beginnt, „damit der Anstieg besser zu sehen ist“. Ist das nötig und fair? Begründe mit deinem Ergebnis aus b).
\end{teile}
\end{aufgabe}
